\documentclass[12pt]{article}

\usepackage[spanish]{babel}
\usepackage[utf8]{inputenc}
\usepackage[T1]{fontenc}
\usepackage{geometry}
\usepackage{array}
\usepackage{longtable}
\usepackage{hyperref}

\geometry{letterpaper,margin=2.5cm}

\begin{document}
	
	
	% Encabezado institucional
	\begin{center}
		{\bfseries UNIVERSIDAD AUTÓNOMA METROPOLITANA}\\[4pt]
		{\bfseries UNIDAD CUAJIMALPA}\\[4pt]
		DIVISIÓN DE CIENCIAS DE LA COMUNICACIÓN Y DISEÑO\\[4pt]
		{\bfseries LICENCIATURA EN TECNOLOGÍAS DE LA INFORMACIÓN}
	\end{center}
	
	\vspace{1em}
	
	\begin{table}[h]
		\renewcommand{\arraystretch}{1.2}
		\begin{tabular}{@{}l@{\hspace{0.5cm}}l@{}}
			PROFESOR: &
			\underline{\makebox[0.7\textwidth][c]{\MakeUppercase{Juan Manuel Alvarado López}}} \\[0.3em]
			ADSCRITO AL: &
			{\small \underline{\makebox[0.7\textwidth][l]{\textbf{DEPARTAMENTO DE TECNOLOGÍAS DE LA INFORMACIÓN}}}} \\[0.3em]
			OTRO: &
			\underline{\makebox[0.7\textwidth][c]{ }} \\[0.3em]
			UEA: &
			\underline{\makebox[0.7\textwidth][c]{\MakeUppercase{Bases de datos avanzadas}}} \\[0.3em]
		\end{tabular}
	\end{table}
	
	
	% Objetivo
	\noindent{OBJETIVO:} El curso aborda conceptos y tecnologías de \textbf{bases de datos avanzadas} más allá del modelo relacional clásico, incluyendo modelos semánticos, bases de datos orientadas a objetos, distribuidas, semiestructuradas y NoSQL, así como aplicaciones analíticas y de grafos. Se busca que el estudiante comprenda las limitaciones del enfoque relacional tradicional y conozca alternativas de modelado, almacenamiento y consulta adecuadas a escenarios complejos, distribuidos o de gran escala.
	
	Al final del curso los alumnos serán capaces de:
	\begin{enumerate}
		\item Conocer las características de la evolución experimentada por la tecnología de gestión de datos.
		\item Emplear los principales modelos semánticos de datos.
		\item Conocer y aplicar los fundamentos de las bases de datos orientadas a objetos.
		\item Conocer la relación entre bases de datos relacionales y orientadas a objetos.
		\item Entender y aplicar los conceptos básicos de organización y diseño de las bases de datos distribuidas.
		\item Explorar las bases de datos semi-estructuradas como alternativa de la organización de las bases de datos.
	\end{enumerate}
	
	\vspace{1em}
	
	% HORARIO en tabla
	\section*{Horario}
	
	\begin{center}
		\renewcommand{\arraystretch}{1.3}
		\begin{tabular}{|c|c|c|}
			\hline
			\textbf{Día} & \textbf{Horario} & \textbf{Aula} \\
			\hline
			\underline{\makebox[2.8cm][c]{Martes}} &
			\underline{\makebox[2.8cm][c]{11--14}} &
			\underline{\makebox[2.8cm][c]{A-608COM}} \\
			\hline
			\underline{\makebox[2.8cm][c]{Viernes}} &
			\underline{\makebox[2.8cm][c]{9--11}} &
			\underline{\makebox[2.8cm][c]{A-428}} \\
			\hline
			%			\underline{\makebox[2.8cm][c]{\ }} &
			%			\underline{\makebox[2.8cm][c]{\ }} &
			%			\underline{\makebox[2.8cm][c]{\ }} \\
			%			\hline
		\end{tabular}
	\end{center}
	
	
	\vspace{1em}


% Plan de 11 semanas (usando tu contenido, resumido en la tabla)
\section*{Plan de trabajo (11 semanas)}

\setlength{\extrarowheight}{2pt}
\begin{longtable}{|>{\centering\arraybackslash}p{1.2cm}|p{7.2cm}|p{7.2cm}|}
	\hline
	\textbf{Semana} & \textbf{Temas principales} & \textbf{Actividades} \\
	\hline
	\endfirsthead
	\hline
	\textbf{Semana} & \textbf{Temas principales} & \textbf{Actividades} \\
	\hline
	\endhead
	
	1 &
	Presentación de la UEA “Bases de Datos Avanzadas”: objetivos general y específicos, métodos de evaluación (exámenes, cuestionarios, prácticas, proyecto final), reglamento y visión general del contenido. & \\ 
	\hline
	
	2 &
	Evolución de la tecnología y limitaciones del modelo relacional (parte 1): recapitulación de modelo relacional clásico (tablas, claves, restricciones, SQL básico) y evolución histórica de archivos planos a SGBD y al modelo relacional. &
	Reflexión sobre nuevas necesidades (datos complejos, multimedia, distribución geográfica, escalabilidad, web). \\
	\hline
	
	3 &
	Evolución de la tecnología y limitaciones del modelo relacional (parte 2). Limitaciones del modelo relacional y primera introducción a modelos semánticos de datos (EER y noción de modelo semántico). &
	Análisis de esquemas relacionales complicados y actividad en equipos para proponer mejoras usando ideas de modelos semánticos. \\
	\hline
	
	4 &
	Modelos semánticos de datos: modelo Entidad–Relación extendido (especialización/generalización, agregación, atributos multivaluados, entidades débiles) y vínculos con modelos orientados a objetos y objeto–relacionales. &
	Comparar un diagrama ER simple con su versión extendida para un mismo problema; Usando un caso de estudio, identificar requisitos difíciles de expresar en un modelo relacional plano y proponer alternativas semánticas. \\
	\hline
	
	5 &
	Bases de datos orientadas a objetos: fundamentos y relación con modelos semánticos. Recordatorio de POO y nociones de objetos persistentes, identificadores de objeto y clases almacenadas en la base de datos. &
	Construcción de un modelo semántico con generalización, composición y atributos complejos. \\
	\hline
	
	6 &
	Examen parcial que cubre los temas vistos hasta el momento & \\
	\hline
	
	7 &
	Bases de datos distribuidas: fundamentos, organización y diseño. Concepto de base de datos distribuida, arquitecturas (centralizadas, distribuidas, federadas), fragmentación (horizontal, vertical, híbrida) y replicación. &
	Transformar parte de un modelo semántico en un diseño orientado a objetos u objeto–relacional y discutir cómo podría distribuirse; proponer estrategias de fragmentación y replicación para un sistema con sedes en varias ciudades. \\
	\hline
	
	8 &
	Bases de datos semiestructuradas: datos con esquema flexible, XML y JSON, bases documentales como complemento a bases relacionales y consulta básica (XPath/XQuery, consultas sobre JSON). &
	Representar información en XML y JSON para un caso de estudio y plantear consultas de extracción y búsqueda. \\
	\hline
	
	9 &
	Otras tecnologías: NoSQL (documentales, clave–valor, columnas y grafos) y aplicaciones de bases de datos avanzadas; motivaciones de escalabilidad, alta disponibilidad y flexibilidad de esquemas. &
	Representar un mismo conjunto de datos en XML y JSON y diseñar consultas; analizar qué tipo de base NoSQL (documental, clave–valor, grafo o columna) conviene para distintos escenarios; uso sugerido de MongoDB Community Edition. \\
	\hline
	
	10 &
	Proyecto final y cierre: asesoría para el proyecto en equipo, revisión de la arquitectura de bases de datos avanzadas elegida. & \\
	\hline
	
	11 &
	Examen final. & \\
	\hline
	
\end{longtable}

	
	% Criterios de evaluación (desde el .md)
	\section*{Criterios e instrumentos de evaluación}
	
	La calificación final se integra a partir de los siguientes instrumentos de evaluación. Cada uno mide distintas dimensiones del aprendizaje: la comprensión conceptual, la aplicación práctica y el desempeño individual.
	
	\begin{center}
		\begin{tabular}{|l|p{7cm}|c|}
			\hline
			\textbf{Instrumento} & \textbf{Descripción} & \textbf{Porcentaje} \\
			\hline
			Proyecto & Trabajo colaborativo integrador de conceptos vistos en clase & 25\% \\
			\hline
			Actividades & Ejercicios, tareas, prácticas, participación & 25\% \\
			\hline
			Examen parcial & Evaluación escrita o práctica a mitad del periodo & 25\% \\
			\hline
			Examen final & Evaluación escrita o práctica al cierre del periodo & 25\% \\
			\hline
			Total & & 100\% \\
			\hline
		\end{tabular}
	\end{center}
	
	\vspace{0.5em}
	
	\noindent\textbf{Calificación mínima aprobatoria.} La calificación mínima aprobatoria es \textbf{6.0 (seis)} en escala de 0 a 10.
	
	\noindent\textbf{Equivalencias numéricas:}
	\begin{itemize}
		\item NA: 0.0 a < 6.0
		\item S: 6.0 a < 8.0
		\item B: 8.0 a < 9.0
		\item MB: 9.0 a 10.0
	\end{itemize}
	
	\subsection*{Entrega de trabajos y actividades}
	
	\begin{itemize}
		\item Las fechas de entrega se publicarán con anticipación en la plataforma institucional y no serán modificadas salvo causa de fuerza mayor debidamente justificada.
		\item Los trabajos se entregarán \textbf{exclusivamente en los formatos indicados} (PDF, Word, ZIP u otro especificado) y a través del canal oficial del curso (por ejemplo, la plataforma LMS institucional).
		\item \textbf{No se aceptarán entregas mediante ligas o vínculos a servicios de almacenamiento en la nube} (Google Drive, OneDrive, Dropbox, etcétera).
		\item Los trabajos entregados después de la fecha y hora límite no serán considerados para evaluación, salvo circunstancias excepcionales documentadas.
	\end{itemize}
	
	% Política de asistencia (desde el .md)
	\section*{Política de asistencia}
	
	La asistencia regular es un factor determinante en el aprendizaje. El seguimiento de las actividades presenciales, la participación en clase y la interacción entre pares no pueden reemplazarse por otros medios.
	
	\subsection*{Requisito mínimo de asistencia}
	
	\textbf{Se espera una asistencia regular a las sesiones del curso.} Cuando la asistencia mínima establecida no se cumple, el estudiante deja de ser elegible para consideraciones especiales en entregas extemporáneas, reprogramación de evaluaciones u otros ajustes extraordinarios.
	
	\subsection*{Definición de asistencia, retardo y falta}
	
	\begin{center}
		\begin{tabular}{|l|p{8cm}|}
			\hline
			\textbf{Condición} & \textbf{Criterio} \\
			\hline
			Asistencia & Presente dentro de los primeros 15 minutos de iniciada la sesión \\
			\hline
			Retardo & Presente entre el minuto 16 y el minuto 30 de iniciada la sesión \\
			\hline
			Falta & Presente después del minuto 30, o no se presenta \\
			\hline
			Falta por abandono & Salida del aula sin justificación por más de 20 minutos continuos durante la sesión \\
			\hline
		\end{tabular}
	\end{center}
	
	\vspace{0.5em}
	
	\noindent 2 retardos equivalen a 1 falta. Se permiten hasta \textbf{3 faltas}. A partir de la cuarta falta, el estudiante pierde el derecho a consideraciones especiales.
	
	\subsection*{Faltas justificadas}
	
	Las inasistencias por motivos de salud, compromisos institucionales (representación deportiva, cultural o académica) o causas de fuerza mayor deberán notificarse \textbf{antes de la sesión o a más tardar 48 horas después}, acompañadas del documento que las acredite. \textbf{Sin documento de evidencia no se justifica la falta.}
	
	% Integridad académica (desde el .md)
	\section*{Integridad académica: plagio y uso de IA}
	
	La formación en este curso busca desarrollar competencias auténticas. La elaboración propia de los trabajos y evaluaciones es parte fundamental de ese proceso.
	
	\subsection*{Plagio}
	
	Se considera plagio toda reproducción total o parcial de ideas, textos, código fuente, diagramas u otros contenidos de terceros, sin el reconocimiento explícito de la fuente original mediante cita adecuada.
	
	\begin{itemize}
		\item El estudiante es responsable de citar correctamente cualquier fuente consultada, ya sea impresa, digital o generada por herramientas de terceros.
		\item Los trabajos que presenten indicios de plagio serán reportados y \textbf{no serán tomados en cuenta para efectos de evaluación}, quedando con valor de cero en el instrumento correspondiente.
	\end{itemize}
	
	\subsection*{Uso de herramientas de inteligencia artificial}
	
	El uso de herramientas de inteligencia artificial generativa (como ChatGPT, Copilot, Gemini u otras similares) está \textbf{sujeto a las indicaciones específicas de cada actividad}. En términos generales:
	
	\begin{itemize}
		\item \textbf{Uso permitido y declarado:} Las herramientas de IA podrán utilizarse como apoyo para exploración, estructuración o revisión de ideas, siempre que el estudiante declare explícitamente su uso al momento de la entrega.
		\item \textbf{Uso no permitido:} Queda prohibido presentar como producción propia contenido íntegramente generado por una IA, sin intervención analítica o creativa sustancial del estudiante.
		\item \textbf{En exámenes:} Queda estrictamente prohibido el uso de cualquier herramienta de IA, navegador web u otro recurso externo no autorizado explícitamente por el docente.
	\end{itemize}
	
	\subsection*{Detección y consecuencias}
	
	Se hará uso de las \textbf{herramientas de detección de plagio y contenido generado por IA} disponibles. Si se detecta plagio o uso indebido de IA en un trabajo o actividad evaluable, dicho instrumento \textbf{quedará sin efecto para efectos de calificación} y se registrará con valor de cero, sin posibilidad de re-entrega. En casos reincidentes, la situación será reportada a las instancias académicas correspondientes. El uso indebido de dispositivos electrónicos durante exámenes también anula la calificación del examen y será reportado.
	
	% Conducta y convivencia (desde el .md)
	\section*{Conducta y convivencia en el aula}
	
	Un ambiente de aprendizaje respetuoso beneficia a todos. Las siguientes normas de convivencia aplican durante todas las sesiones presenciales:
	
	\begin{itemize}
		\item \textbf{Puntualidad:} Se espera la presencia del estudiante al inicio de cada sesión.
		\item \textbf{Respeto:} Se mantiene un trato respetuoso hacia los compañeros de clase y el personal de la institución en todo momento.
		\item \textbf{Dispositivos electrónicos:} El uso de computadoras, tabletas y teléfonos inteligentes está \textbf{permitido únicamente con fines académicos} relacionados con la sesión en curso.
		\item \textbf{Silencio y atención:} Se solicita silenciar notificaciones y evitar conversaciones paralelas que interrumpan la dinámica de la clase.
		\item \textbf{Cuidado de instalaciones:} El estudiante es responsable de mantener en buen estado el mobiliario, equipo y espacios del aula.
	\end{itemize}
	
	
	
\end{document}